% !Mode:: "TeX:UTF-8"
\chapter{系统详细设计与主要算法}

本章依据第三章给出的系统总体架构图展开详细设计，
重点描述算法处理层中的核心模块，
以及这些模块与数据层之间的必要关联，
包括数据接入与特征表示、异构存储、故障辅助分析以及数据分析与优化等内容。

\section{数据接入与特征表示设计}

数据接入与特征表示设计对应架构图中的数据接入与管理模块。
该部分负责把Fuzzing产物中的日志、测试脚本和算子拓扑结构整理为可检索、可分析的信息。
其中，日志解析用于提取错误类型和运行状态，
双路向量化用于把日志语义和模型结构分别映射到向量空间。

\subsection{基于分层正则匹配的日志解析模块}

Fuzzing工具输出的崩溃日志在格式上不统一。
同一条记录中可能同时包含多层运行时信息以及多线程交织产生的乱序输出。
日志解析模块接收原始崩溃日志及相关测试信息，
将其中稳定的字段提取为后续分析所需的元数据。

本系统设计的LogParser模块采用分层正则匹配策略，
将解析过程分解为错误分类与字段提取两个独立阶段。
在错误分类阶段，模块维护一张带有优先级顺序的错误模式表，
涵盖八种PyTorch测试中常见的错误类别，如表\ref{tab:error_categories}所示。

\begin{table}[htbp]
\centering
\caption{LogParser模块支持的八种标准化错误类别}
\label{tab:error_categories}
\begin{tabular}{lll}
\toprule
类别标识 & 中文描述 & 典型触发关键词 \\
\midrule
\texttt{cuda\_oom} & CUDA显存溢出 & cuda out of memory, oom \\
\texttt{shape\_mismatch} & 张量维度不匹配 & size mismatch, incompatible \\
\texttt{segfault} & 内存段错误 & segmentation fault, signal 11 \\
\texttt{timeout} & 执行超时 & timeout, exceeded maximum \\
\texttt{device\_mismatch} & 设备不匹配 & expected.*device \\
\texttt{index\_error} & 索引越界 & index out of, out of bounds \\
\texttt{type\_error} & 类型错误 & type error, invalid type \\
\texttt{nan\_inf} & 数值异常 & nan, inf, not a number \\
\bottomrule
\end{tabular}
\end{table}

模块将输入文本统一转换为小写后，依照表中自上而下的优先级顺序扫描关键词。
一旦命中即终止扫描并返回对应错误类别。
优先级排列主要依据异常处理时的排查优先级：
显存溢出与段错误属于系统级致命异常，排在最前；
数值异常属于静默型偏差，排在最后。
对于不命中任何模式的日志，模块将其归入\texttt{other}类别。
该分类结果用于为后续故障辅助定位和检索分析提供错误标签。

字段提取阶段则针对不同数据类型分别准备了一组正则表达式，
从原始日志中抽取出以下六项结构化字段：

（1）错误类型（error\_type）。
面向Python异常，模块使用形如\texttt{(\textbackslash w+Error):}的正则匹配得到类名；
遇到C++层的段错误，则改为检测是否出现\texttt{Segmentation Fault}关键字。

（2）张量维度（tensor\_shapes）。
PyTorch的维度报错主要包含方阵表示法\texttt{m1: [32 x 512]}，
部分报错还会直接给出匿名形状\texttt{[32 x 512]}。
模块为这两类格式分别编写正则表达式并按优先级依次尝试，
成功匹配后输出为\texttt{\{name, shape\}}结构化字典列表。

（3）框架源码位置（source\_location）。
C++侧报错一般带有类似\texttt{at /pytorch/.../}\allowbreak\texttt{THTensorMath.cpp:41}的源码路径，
Python侧则使用\texttt{File "xxx.py", line 42}这样的格式。
模块分别准备了正则表达式，用于取出文件名、行号和模块路径。

（4）显存信息（memory\_info）。
仅在\texttt{cuda\_oom}类别下激活。
模块提取三组数值：请求分配量、已分配量与设备总容量。

（5）设备信息（device\_info）。
通过全局扫描\texttt{cuda:N}与\texttt{cpu}标识符提取，去重后输出为列表。

（6）原始消息（raw\_message）。
保留未经处理的完整日志文本，供RAG模块作为LLM上下文素材。

\subsection{双路向量化特征表示方法}

如第二章所述，本系统采用Sentence-BERT编码模型将离散文本映射为384维稠密浮点向量。
针对Fuzzing产物的异构特性，系统设计了两条并行的向量化路径。

第一条路径为报错堆栈语义表征。
模块把用例的\texttt{error\_message}字段直接送入Sentence-BERT编码器。
遇到成功执行、没有错误信息的用例，则补充占位文本
\texttt{"Success: [status] - Model executed OK"}进行编码。
这条路径产出的向量写入Milvus中的\texttt{fuzzing\_logs}集合，
供后续语义检索使用。

第二条路径为图拓扑序列投影。
模块将测试用例的模型结构字典序列化为标准化文本，
格式为\texttt{"Model structure: layers=[...],}\allowbreak\texttt{ connections=...,}\allowbreak\texttt{ depth=N,}\allowbreak\texttt{ operators=[...]"}。
该文本经Sentence-BERT编码后产出的向量存入Milvus的\texttt{fuzzing\_structure}集合，
用于多样性分析与拓扑相似性召回。

为了避免重复加载开销，两条路径共用同一个all-MiniLM-L6-v2模型实例。
每个测试用例经过该处理后会得到两个独立的384维向量。

\section{异构存储模块设计}

异构存储模块用于连接数据接入的结果与后续分析算法。
结合第三章的技术栈选型，
系统采用MySQL与Milvus组合保存测试用例数据。

在数据组织上，MySQL中的\texttt{test\_case\_metadata}表以\texttt{case\_uid}作为业务标识，
保存用例状态、错误信息、分支边覆盖数、执行耗时和扩展参数等信息。
Milvus维护\texttt{fuzzing\_logs}和\texttt{fuzzing\_structure}两个集合，
分别对应报错堆栈语义向量和模型结构向量。
两类数据不放在同一数据库中，
主要是因为结构化筛选和相似度检索对存储模型的要求不同。

跨库关联通过MySQL中保存的向量主键完成。
数据接入阶段，系统先完成日志文本和结构文本的向量化，
再将两路向量写入Milvus并取得各自的主键标识
（\texttt{structure\_vector\_id}与\texttt{log\_vector\_id}）。
结构向量主键写入MySQL的\texttt{vector\_id}字段，
结构向量主键与日志向量主键同时写入\texttt{parameters}字段，
并与结构化元数据和\texttt{case\_uid}一同保存。
后续检索时，系统可先在Milvus中召回相似向量，
再依据返回的主键查询MySQL中的元数据。
这种设计避免在Milvus中保存业务字段，
也便于在MySQL中维护用例状态、覆盖数据和清退标记。

\section{故障辅助分析模块设计}

故障辅助分析模块对应系统架构图中的故障辅助分析能力，
主要包括单用例故障辅助定位、语义相似召回、故障关联分析和RAG辅助诊断。
其中，故障辅助定位用于在用例详情中标注候选算子节点；
语义相似召回和RAG辅助诊断用于组织历史案例并生成辅助诊断报告；
故障关联分析则结合相似成功样本与崩溃样本，
提取候选边界条件和参数差分。

\subsection{基于DFS拓扑回溯的故障辅助定位模块}

日志解析完成后，系统需要将结构化错误信息与测试用例的计算图结构相结合，
定位出可能与崩溃相关的候选算子节点位置。
本系统设计的FaultLocator模块通过逆向邻接表与DFS回溯实现该功能，
帮助用户缩小排查范围。

逆向邻接表的构建过程如下。
系统将测试用例中的算子拓扑统一表示为有向图$G=(V,\mathcal{E})$，
其中节点集合$V$表示算子节点，边集合$\mathcal{E}$表示数据的流向。
设正向边$(o_i,o_j)\in\mathcal{E}$表示数据由$o_i$流向$o_j$，
则逆向邻接表中记录$o_j$到$o_i$的反向关系。

具体定位过程如算法\ref{alg:dfs_fault_locator}所示。

\begin{algorithm}[htbp]
\caption{基于DFS回溯的候选算子定位算法}
\label{alg:dfs_fault_locator}
\begin{algorithmic}[1]
\REQUIRE 结构化错误信息$E$，算子有向图$G=(V,\mathcal{E})$
\ENSURE 候选算子节点列表$\mathcal{L}$
\STATE 根据边集合$\mathcal{E}$构建逆向邻接表$\mathbf{A}^{-1}$
\STATE 根据$E$中的错误类别选择错误相关算子集合$\mathcal{M}$
\STATE 根据$E$确定候选触发节点集合$\mathcal{T}$；若无法确定，则令$\mathcal{T}\leftarrow V$
\STATE $\mathcal{L} \leftarrow \emptyset$
\FOR{each trigger node $t \in \mathcal{T}$}
    \STATE 从$t$开始沿$\mathbf{A}^{-1}$执行DFS回溯
    \FOR{each visited node $o_i$}
        \IF{$o_i \in \mathcal{M}$ or $o_i$ matches a risky topology pattern}
            \STATE 计算候选置信度$score(o_i)$并生成推理理由
            \STATE $\mathcal{L} \leftarrow \mathcal{L} \cup \{(o_i, \mathrm{id}(o_i), score(o_i), reason)\}$
        \ENDIF
        \IF{$o_i \in \mathcal{M}$ or depth exceeds threshold or $\mathbf{A}^{-1}(o_i)=\emptyset$}
            \STATE 停止当前回溯分支
        \ENDIF
    \ENDFOR
\ENDFOR
\STATE 对$\mathcal{L}$按置信度降序排序
\RETURN $\mathcal{L}$
\end{algorithmic}
\end{algorithm}

其中，错误相关算子集合$\mathcal{M}$根据错误类型确定。
例如，维度变换错误主要关注Linear、Conv2d和Reshape等算子；
高显存问题优先检查Conv类和Attention类算子。
这里的置信度只用于排序候选节点，不表示真实缺陷概率。
该列表将传递给前端在计算图可视化中高亮标注候选节点。

\subsection{语义相似召回与RAG辅助诊断}

当用户针对某一崩溃用例发起诊断请求时，
系统先从历史用例中召回语义相近的记录，
再把召回结果作为RAG上下文交给LLM生成辅助诊断报告。
其处理过程如算法\ref{alg:rag_retrieval}所示。

\begin{algorithm}[htbp]
\caption{语义相似召回与RAG辅助诊断算法}
\label{alg:rag_retrieval}
\begin{algorithmic}[1]
\REQUIRE 用户查询$q$，最终返回数量$K$
\ENSURE 召回用例集合$\mathcal{C}$，辅助诊断报告$R$
\STATE $v_q \leftarrow \text{SentenceBERT}(q)$
\STATE 在\texttt{fuzzing\_structure}中检索与$v_q$最相似的前$3K$个结构候选
\STATE 在\texttt{fuzzing\_logs}中检索与$v_q$最相似的前$3K$个日志候选
\STATE 根据$q$中的结构关键词和日志关键词比例计算两路融合权重
\STATE 合并两路候选结果，并按融合得分截取前$K$个用例，记为$\mathcal{C}$
\STATE 根据$\mathcal{C}$中的主键从MySQL反查完整元数据
\STATE 将元数据和用户查询$q$组装为RAG上下文提示词$P$
\STATE $R \leftarrow \text{LLM}(P)$
\RETURN $\mathcal{C}, R$
\end{algorithmic}
\end{algorithm}

LLM基于检索到的真实历史案例生成辅助诊断报告，
为开发者提供故障排查参考信息。

\subsection{故障关联分析}

故障关联分析中的候选边界条件分析模块借鉴经典的Daikon动态不变量推断范式\cite{ref_daikon}，
在成功样本与崩溃样本构成的对照组中扫描候选约束，
输出候选边界条件，用于辅助分析崩溃触发条件。

候选边界条件分析的输入为一个目标崩溃用例。
模块先提取该用例的图拓扑结构向量，
在Milvus的\texttt{fuzzing\_structure}集合中执行余弦相似度检索，
分别召回拓扑结构相似但执行成功的10个用例（健康组$\mathcal{S}$）
与执行失败的10个用例（崩溃组$\mathcal{F}$）。
召回条件为结构相似度大于0.85，确保对照样本与目标用例的算子组合足够相近。

召回完成后，模块对每个用例的嵌套参数字典执行展平操作。
DL测试用例的参数空间通常包含模型结构字典、层配置子字典、标量参数和列表参数。
展平操作把这些嵌套结构转成一维键值对，
同时过滤向量主键、覆盖率等无关字段。
过滤后的键值对集合仅保留与算子配置直接相关的参数。

展平完成后，模块计算健康组与崩溃组所有样本的键集合交集，
其中$\Phi(\cdot)$表示参数展平后的键值对集合：
\begin{equation}
\label{eq:common_keys}
\mathcal{K} = \bigcap_{s \in \mathcal{S}} \text{keys}(\Phi(s)) \;\cap\; \bigcap_{f \in \mathcal{F}} \text{keys}(\Phi(f))
\end{equation}
仅在$\mathcal{K}$中的键才进入后续的候选约束扫描环节。

模块在公共键集合$\mathcal{K}$上扫描三层模板池，
包括单变量阈值与奇偶性模板、双变量约束与整除关系模板，
以及三变量约束模板。
仅当某条候选约束同时满足“成功全票通过”与“失败全票否决”的双重验证时，
该约束才被加入候选边界条件集合。

上述三层扫描的双重验证判据以算法\ref{alg:dual_verify}的伪代码形式给出。

\begin{algorithm}[htbp]
\caption{双重验证判据——候选边界条件提取算法}
\label{alg:dual_verify}
\begin{algorithmic}[1]
\REQUIRE 健康组展平参数集合 $\mathcal{S}_{\text{flat}} = \{s_1, s_2, \ldots, s_{|\mathcal{S}|}\}$
\REQUIRE 崩溃组展平参数集合 $\mathcal{F}_{\text{flat}} = \{f_1, f_2, \ldots, f_{|\mathcal{F}|}\}$
\REQUIRE 公共键集合 $\mathcal{K} = \{k_1, k_2, \ldots, k_n\}$
\ENSURE 候选约束集合 $\mathcal{R} = \{\}$
\FOR{each candidate predicate $\phi(k_{i_1}, \ldots, k_{i_p})$ in template pool}
    \STATE $\textit{success\_pass} \leftarrow \text{True}$
    \STATE $\textit{failure\_reject} \leftarrow \text{True}$
    \FOR{each $s \in \mathcal{S}_{\text{flat}}$}
        \IF{$\phi(s[k_{i_1}], \ldots, s[k_{i_p}])$ is False}
            \STATE $\textit{success\_pass} \leftarrow \text{False}$
            \STATE \textbf{break}
        \ENDIF
    \ENDFOR
    \IF{$\textit{success\_pass}$ is True}
        \FOR{each $f \in \mathcal{F}_{\text{flat}}$}
            \IF{$\phi(f[k_{i_1}], \ldots, f[k_{i_p}])$ is True}
                \STATE $\textit{failure\_reject} \leftarrow \text{False}$
                \STATE \textbf{break}
            \ENDIF
        \ENDFOR
    \ENDIF
    \IF{$\textit{success\_pass}$ is True \AND $\textit{failure\_reject}$ is True}
        \STATE $\mathcal{R} \leftarrow \mathcal{R} \cup \{\phi\}$
    \ENDIF
\ENDFOR
\RETURN $\mathcal{R}$
\end{algorithmic}
\end{algorithm}

提取出的约束集合$\mathcal{R}$经去重后按优先级排序：
多变量约束排在单变量约束之前，
因为多变量约束包含更丰富的变量间交互信息。
排序后的约束列表即为该崩溃用例的候选边界条件集合，
用于为开发者提供后续排查线索。

最后，系统把候选约束集合$\mathcal{R}$、参数差分和相关测试上下文组装为提示词，
由LLM生成用于解释候选约束的辅助诊断报告。
前端将候选数学公式和辅助诊断报告并列展示。

\section{数据分析与优化模块设计}

数据分析与优化模块对应架构图中的数据分析与优化能力。
该模块针对测试集整体质量，
提供多样性分布观察、效能评估、算子覆盖缺口分析以及候选补盲测试脚本生成等功能。

\subsection{基于PCA的测试集多样性分析}

Fuzzing过程会不断产生测试用例。
这些用例在语义向量空间中的分布，
可以用于观察测试集在算子组合和参数结构上的差异。
由于384维向量空间不便直接查看，
系统需要先将它压缩到二维平面。

本系统在批量更新后或用户请求多样性分析时，对当前样本集合统一执行PCA降维。
具体流程为：取当前参与分析的测试用例结构向量矩阵$\mathbf{X} \in \mathbb{R}^{b \times 384}$
（$b$为参与降维的样本数量），计算协方差矩阵并提取前两个主成分方向，
将每个样本投影至二维坐标$(x_i, y_i)$。
投影后的坐标经最小-最大归一化映射至$[-1, 1]$区间，
存入MySQL的\texttt{parameters}字段中的\texttt{vis\_x}与\texttt{vis\_y}键。

前端散点图基于上述二维坐标渲染。
用户可观察测试集的空间分布形态：
若散点呈现明显的簇状聚集，说明当前Fuzzing策略可能集中在有限的算子组合区域内；
若散点分布较为分散，说明测试输入在模型结构或参数组合上具有更高差异。
稀疏区域对应尚未被充分测试的算子拓扑组合，
可作为后续补盲策略的参考方向。

\subsection{CEI覆盖效能指数量化模型}

CEI用来描述单个测试用例在覆盖贡献数据和计算成本之间的关系，其数学定义如下：
\begin{equation}
\label{eq:cei}
\mathrm{CEI} = \frac{E_{\text{edge}}}{\displaystyle\sum_{i=1}^{m} w_i \cdot c_i}
\end{equation}
其中$E_{\text{edge}}$为该用例的分支边覆盖数，
$m$为该用例计算图中不同算子类型的数量，
$c_i$为第$i$类算子在计算图中的出现次数，
$w_i$为第$i$类算子对应的权重系数。
公式分母表示按算子类型和出现次数估计的计算成本，
不包含执行耗时变量；执行耗时可以作为实验评估指标单独记录。

算子权重$w_i$的设计原则是：
计算量越大、参数空间越大的算子，赋予的权重越大，
因为这类算子消耗的执行资源更多，却不一定带来更高的覆盖贡献。
表\ref{tab:operator_weight}列出本系统的部分算子权重。

\begin{table}[htbp]
\centering
\caption{部分算子的CEI权重基准}
\label{tab:operator_weight}
\begin{tabular}{llr}
\toprule
算子类别 & 典型算子 & 权重$w_i$ \\
\midrule
计算密集型 & Conv2d, Conv3d, Linear, LSTM & 3.0 \\
注意力机制 & MultiHeadAttention, Transformer & 4.0 \\
归一化层 & BatchNorm2d, LayerNorm, GroupNorm & 1.5 \\
池化层 & MaxPool2d, AvgPool2d, AdaptiveAvgPool & 1.0 \\
激活函数 & ReLU, Sigmoid, Tanh, GELU & 0.5 \\
结构操作 & Flatten, Reshape, Dropout, Concat & 0.3 \\
\bottomrule
\end{tabular}
\end{table}

系统在执行清退时，把全库CEI中位数的$\alpha$倍作为淘汰阈值。
CEI低于这个阈值的用例在当前指标下覆盖收益相对较低，会进入级联删除流程。
这一流程主要帮助识别冗余数据，
减少它们对存储空间和后续分析结果的干扰。

当用户确认执行清退操作时，
系统需要按照MySQL中保存的向量主键同步处理Milvus向量记录。
具体流程分为三步。
第一步，从MySQL中批量查询满足清退条件的记录，
提取其\texttt{vector\_id}、
\texttt{parameters.structure\_}\allowbreak\texttt{vector\_id}
和\texttt{parameters.log\_vector\_id}。
第二步，调用Milvus的\texttt{delete}接口按主键删除两个集合中的对应向量。
第三步，在MySQL中执行\texttt{DELETE}语句移除对应元数据记录。

\subsection{算子覆盖缺口分析}

除了对单个用例做效能评估之外，系统还需要从全局视角分析测试盲区，
即在当前已出现算子和常用算子范围内，
尚未被测试集覆盖的算子组合模式。

系统不直接枚举PyTorch中的全部公开算子，
而是构造用于缺口分析的候选算子集合$\mathcal{O}$。
该集合由两部分组成：
一部分是当前测试集中已经出现过的算子类型，
另一部分是预设的20个常用PyTorch算子类型。
记当前测试集已出现算子集合为$\mathcal{O}_{\text{seen}}$，
常用算子集合为$\mathcal{O}_{\text{common}}$，
则有：
\begin{equation}
\label{eq:operator_pool}
\mathcal{O} = \mathcal{O}_{\text{seen}} \cup \mathcal{O}_{\text{common}}
\end{equation}

全局邻接矩阵$\mathbf{A} \in \{0,1\}^{|\mathcal{O}| \times |\mathcal{O}|}$
记录了当前测试集中已经出现过的算子对组合关系。

盲区提取通过集合差运算来完成。
定义候选可行算子对全集$\mathcal{E}_{\text{all}}$为$\mathcal{O}$中
所有在PyTorch算子语义下合法的算子组合（排除不合理的组合，
例如\texttt{Softmax}不能直接接\texttt{Conv2d}的通道输入端）。
已覆盖算子对集合$\mathcal{E}_{\text{covered}} = \{(i,j) \mid \mathbf{A}[i][j] = 1\}$。
盲区集合定义为：
\begin{equation}
\label{eq:blind_spots}
\mathcal{B} = \mathcal{E}_{\text{all}} \setminus \mathcal{E}_{\text{covered}}
\end{equation}

盲区集合$\mathcal{B}$中的每个元素$(o_i, o_j)$表示一对还没被测试的算子衔接。
系统把$\mathcal{B}$按算子类别分组后输出为结构化列表。
这一列表会作为后续补盲模块的输入目标。

\subsection{覆盖缺口驱动的候选补盲测试脚本生成}

盲区集合$\mathcal{B}$标记了当前测试集的覆盖缺口。
本系统通过构造结构化提示词，
让LLM生成候选补盲测试脚本来覆盖这些缺口。

结构化提示词的设计如表\ref{tab:prompt_template}所示：

\begin{table}[htbp]
\centering
\caption{RAG补盲提示词模板结构}
\label{tab:prompt_template}
\begin{tabular}{lp{10cm}}
\toprule
板块 & 内容设计 \\
\midrule
\textbf{System} &
“你是一个深度学习模糊测试专家助手。\allowbreak{}
你的任务是根据用户问题和检索到的测试用例上下文，\allowbreak{}
提供专业的分析和建议。请基于提供的测试用例数据给出具体建议；\allowbreak{}
如果问题涉及错误分析，给出可疑条件和排查方向；\allowbreak{}
如果问题涉及优化建议，给出可行的改进方案。\allowbreak{}
保持回答简洁专业，使用中文回答。” \\
\midrule
\textbf{Context} &
注入从Milvus召回的$K$个最相似历史用例的结构化摘要，
每条包含：用例编号、执行状态和错误信息摘要等。
格式为“用例N: [case\_uid] - 状态: [status] -
算子: [operators]”。\\
\midrule
\textbf{Target} &
明确列出当前需要补盲的算子对列表，
例如“请生成一个包含\allowbreak{}\texttt{Conv2d->}\allowbreak\texttt{GroupNorm->LSTM}\allowbreak{}连接链的PyTorch模型脚本”。
包含目标算子的参数范围约束
（如\texttt{in\_channels}$\in [1, 512]$、\allowbreak{}%
\texttt{kernel\_size}$\in [1, 7]$）。\\
\midrule
\textbf{Constraint} &
语法限制与输出格式要求：
“输出应为可检查的候选Python脚本；
模型应继承\texttt{torch.nn.Module}；\allowbreak{}
\texttt{forward}方法应接受形状为\allowbreak{}%
\texttt{[batch, channels, H, W]}的输入张量；
脚本末尾应包含推理调用语句。”\\
\bottomrule
\end{tabular}
\end{table}

上述四板块模板在工程实现中的作用如下：
System板块限定LLM的角色边界，减少它偏离测试生成任务；
Context板块注入真实的历史数据作为参考，减少无上下文生成的问题；
Target板块指定覆盖目标，减少LLM生成已有冗余用例的可能；
Constraint板块给出语法和结构约束，让输出更便于后续检查和执行验证。

LLM生成的候选补盲测试脚本经过必要的检查和确认后，
可以作为后续测试输入的候选材料，
再由开发者确认是否进入数据接入流程，
形成“发现盲区→生成候选补盲测试脚本→执行验证→后续分析”的处理流程。